\color{unidarkgreen}{\subsection{Логика отключения общие}}
\color{black}

\begin{enumerate}[label=\arabic{section}.\arabic{subsection}.\arabic*,  align=left, labelsep=4pt, leftmargin=0pt, itemindent=57pt]

\item
Срабатывание логики отключения основных защит формируются при включенной накладке \textbf{«ЛО
осн»} следующими сигналами:
\begin{itemize}
\item Откл.ДТЗ;
\item Отключение от ГЗ АТ сигн;
\item Ср. ГЗ АТ откл;
\item Ср. ГЗ РПН;
\item Отключение от ГЗ ЛРТ сигн;
\item Ср. ГЗ ЛРТ откл;
\item Ср.отсеч.клап. АТ;
\item Ср.предохр.клап. АТ;
\item Ср.предохр.клап. ЛРТ;
\item Откл. от ТЗ АТ;
\item Откл. от ТЗ ЛРТ;
\item Ср. ЗПО.
\end{itemize}

\item
Действие основных защит на отсечной клапан вводится накладкой \textbf{«Действие на ОК»}, время
импульса действия регулируется \textbf{«Tимп ОК»}.

\item
Срабатывание логики отключения резервных защит формируются при включенной накладке
\textbf{«ЛО рез»} следующими сигналами:
\begin{itemize}
\item Ср. МТЗ-1 НН;
\item Ср. МТЗ-2 НН.
\end{itemize}


\item
Логика деления от МТЗ-1 НН1(НН) выводится накладкой \textbf{«ЛД НН от МТЗ-1 НН»}.
Срабатывание логики деления формируются при возникновении пуска МТЗ-1 НН. Так же
срабатывание контролируется выдержкой времени \textbf{«Тср МТЗ-1 на В НН»}.

\item
Логика деления от МТЗ-2 НН1 выводится накладкой \textbf{«ЛД НН2 от МТЗ-2 НН»}. Срабатывание
логики деления формируются при возникновении пуска МТЗ-2 НН. Так же срабатывание
контролируется выдержкой времени \textbf{«Тср МТЗ-2 на В НН»}.


\begin{figure}[H]
\centering
\includegraphics[width=0.8\textwidth,height=0.8\textheight,keepaspectratio]{\fbpath/041.9902 ЛО общ/img/LO_gen.pdf}
\captionof{figure}{Функциональная схема общей логики отключения}\label{pic:LO_gen}
\end{figure}


\small
\begin{longtable}{|>{\centering\arraybackslash}m{5.3cm}|>{\centering\arraybackslash}m{4.2cm}|>{\centering\arraybackslash}m{1cm}|>{\centering\arraybackslash}m{3.3cm}|>{\centering\arraybackslash}m{1.8cm}|}
% Первая страница: полный заголовок
\caption{Перечень уставочных параметров функции общей логики отключения\hfill\vspace{-0.5\baselineskip}}\label{desc:tbl1}\\
\hline
\rowcolor{gray!30}
Наименование уставки & Диапазон & Шаг & Значение по умолчанию & Единица измерения \\
\hline
\endfirsthead

% Промежуточные страницы: «Продолжение таблицы»
\caption*{\hspace{3pt}\emph{Продолжение таблицы \ref{desc:tbl1}}\hfill\ }\vspace{-0.5\baselineskip} \\
\hline
\rowcolor{gray!30}
Наименование уставки & Диапазон & Шаг & Значение по умолчанию & Единица измерения \\
\hline
\endhead

%===m>ДТЗ|-
%===m
\end{longtable}
\normalsize	

\end{enumerate}